\begin{table}[htbp]\centering
\def\sym#1{\ifmmode^{#1}\else\(^{#1}\)\fi}
\begingroup
\renewcommand{\arraystretch}{1.3}
\caption{Extensive Margin on Agricultural Production, 1970-2010}
\begin{tabular}{l*{4}{c}}
\hline\hline
                    &\multicolumn{2}{c}{Soybean}                &\multicolumn{2}{c}{Cotton}                 \\\cmidrule(lr){2-3}\cmidrule(lr){4-5}
                    &\multicolumn{1}{c}{(1)}         &\multicolumn{1}{c}{(2)}         &\multicolumn{1}{c}{(3)}         &\multicolumn{1}{c}{(4)}         \\
                    &         LPM         &      Probit         &         LPM         &      Probit         \\
\hline              
$\Delta Log Travel Time$&       0.335\sym{*}  &   0.364\sym{**} &       0.095         &       0.108         \\
                    &     (0.180)         &     (0.175)         &     (0.074)         &     (.102)         \\
\hline
Observations        &         528         &         384         &         528         &         381         \\
First-stage F stat  &      58.304         &                     &      58.304         &                     \\
\hline\hline
\end{tabular}
\endgroup
\begin{minipage}{11cm}
    \vspace{0.1cm}
    {\footnotesize \textit{Notes: }  Each observation is a micro-region. Log change in travel time on roads is instrumented by the change in travel times on the predicted road network. The predicted road network is created by geo-referencing the PNV National Transportation Plan in 1973. Dependent variables are defined as a dummy equal to 1 if the micro-region did not produce the crop in 1970 but started to produce after the Plan for National Integration. The reported coefficients for Probit Regressions are the directly interpretable marginal effect. ***,**,* imply significance at the 1\%, 5\%, and 10\% levels, respectively.}
\end{minipage}
\end{table}
